\newcommand{\both}[2]{\ifspanish #2\else #1\fi}
\newcommand{\slsh}[1]{\not\!#1}
\setlength{\emergencystretch}{2em}
\ifspanish\renewcommand{\abstractname}{Resumen}\renewcommand{\figurename}{Figura}\renewcommand{\refname}{Referencias}\fi
\title{\both{Compton production and detection of dark photons}{Producción y detección Compton de fotones oscuros}\\\large TEXONO: \both{vacuum calculation, medium correction, and published limits}{cálculo en vacío, corrección del medio y límites publicados}}
\author{ROOT/C++}\date{25 September 2026}
\begin{document}\maketitle
\begin{abstract}
\both{The tree-level production cross section passes independent matrix, Ward-identity, total-integral and massless-limit checks. Recalculating Park's Figure 1 does not reproduce the full published spectrum; the disagreement is displayed without fitting. For TEXONO we compare Park's quoted bound, the curve extracted from Danilov's PDF, an independent calculation with Park's conventions, and a homogeneous-medium calculation. The medium result includes an explicit signal-selection efficiency. No exclusion is claimed at material resonance or outside the stated small-mixing domain.}
{La sección eficaz de producción a nivel árbol satisface controles matriciales independientes, identidades de Ward, integrales totales y el límite sin masa. El recálculo de la figura 1 de Park no reproduce el espectro publicado completo; mostramos la discrepancia sin ajustar una normalización. Para TEXONO comparamos el límite citado por Park, la curva extraída del PDF de Danilov, un cálculo con las convenciones de Park y un cálculo en medio homogéneo. El resultado en el medio incluye explícitamente la eficiencia de selección. No afirmamos una exclusión en la resonancia material ni fuera del régimen indicado de mezcla pequeña.}
\end{abstract}
\section{\both{What is derived, and what is taken from sources}{Qué se deriva y qué se toma de las fuentes}}
\both{We reuse and recheck the project's v2 polynomial trace kernel. QED vertices, the electron propagator, spin completeness and two-body phase space are input identities. The compact trace, numerical convolution, inverse cross section and exclusion solver are implemented here. Gondolo--Raffelt (Park's Reference 15) supplies an independent total-cross-section check; it is not required to generate the differential kernel. Park supplies the empirical source and experimental inputs. Danilov supplies the medium probabilities and selection correction.}
{Reutilizamos y verificamos de nuevo el núcleo de trazas polinómicas de v2. Los vértices QED, el propagador del electrón, la completitud de espines y el espacio de fases de dos cuerpos son identidades de partida. La traza compacta, la convolución numérica, la sección inversa y el cálculo de exclusión se implementan aquí. Gondolo--Raffelt (referencia 15 de Park) aporta un control independiente de la sección total; no hace falta para generar el núcleo diferencial. Park aporta el espectro empírico y las entradas experimentales. Danilov aporta las probabilidades en el medio y la corrección de selección.}
\section{\both{Production amplitude and trace}{Amplitud de producción y traza}}
\begin{equation}\gamma(k)+e(p)\to X(q)+e(p'),\quad p^2=p'^2=m^2,\ k^2=0,\ q^2=M^2,\quad g={\rm diag}(1,-1,-1,-1).
\end{equation}
\begin{align}
\mathcal L_{\rm int}&=-e\bar\psi\gamma^\mu\psi(A_\mu+\epsilon X_\mu),\\
\mathcal M&=\epsilon e^2\bar u(p')\left[\gamma^\nu\frac{\slsh p+\slsh k+m}{s-m^2}\gamma^\mu+\gamma^\mu\frac{\slsh p-\slsh q+m}{u-m^2}\gamma^\nu\right]u(p)\varepsilon_\mu\xi^*_\nu.
\end{align}
\both{The two electron-exchange diagrams are both required. On-shell Dirac equations give zero upon contraction with either external vector momentum. Therefore the massive polarization sum can be replaced by the metric in the squared amplitude; the longitudinal term vanishes by current conservation. This does not discard a physical polarization.}
{Se requieren ambos diagramas de intercambio del electrón. Las ecuaciones de Dirac sobre la capa de masa anulan la contracción con cualquiera de los momentos vectoriales externos. La suma de polarizaciones masivas puede entonces reemplazarse por la métrica en la amplitud al cuadrado: el término longitudinal se anula por conservación de corriente. Esto no elimina una polarización física.}
\begin{align}
\overline{|\mathcal M|^2}&=\epsilon^2e^4 F,\qquad a=s-m^2,\ b=u-m^2,\ c=m^2,\ d=M^2,\\
F&=-2\left(\frac ab+\frac ba\right)+4(2c+d)\left[\frac1a+\frac1b+c\left(\frac1a+\frac1b\right)^2-\frac d{ab}\right].
\end{align}
\both{The spin trace averages over two electron spins and two incoming photon polarizations, a factor of one quarter. It follows by expanding both diagrams and applying the even gamma-matrix trace identities. The code checks the polynomial identity, then evaluates the same amplitude using independent complex Dirac matrices and explicit physical polarizations.}
{La traza promedia sobre dos espines del electrón y dos polarizaciones del fotón incidente, con un factor de un cuarto. Se obtiene expandiendo ambos diagramas y aplicando las identidades de trazas de matrices gamma pares. El código verifica la identidad polinómica y evalúa la misma amplitud con matrices de Dirac complejas independientes y polarizaciones físicas explícitas.}
\section{\both{Laboratory cross section and threshold}{Sección eficaz en laboratorio y umbral}}
\begin{align}
s&=m^2+2mE_\gamma,\quad b=M^2-2mE_X,\quad E_{\gamma,\rm th}=M+\frac{M^2}{2m},\\
\lambda(s,m^2,M^2)&=[s-(m+M)^2][s-(m-M)^2],\\
E_X^\pm&=\frac{(E_\gamma+m)(s+M^2-m^2)\pm E_\gamma\sqrt{\lambda}}{2s},\\
\frac{\dd\sigma_{\gamma e\to Xe}}{\dd E_X}&=\frac{\epsilon^2e^4(\hbar c)^2}{32\pi mE_\gamma^2}F,\qquad E_X^-\le E_X\le E_X^+.
\end{align}
\both{Outside the endpoints or below threshold the cross section is zero. Integrating the expression numerically agrees with the derived analytic primitive and Reference 15. At zero mass the integrated result is Klein--Nishina times the squared mixing. These are free-electron, tree-level statements; they do not establish that Compton scattering equals total material attenuation.}
{Fuera de los extremos o por debajo del umbral la sección es cero. La integral numérica coincide con la primitiva analítica derivada y con la referencia 15. A masa cero, la integral es Klein--Nishina multiplicada por la mezcla al cuadrado. Son resultados a nivel árbol para electrones libres; no establecen que Compton sea la atenuación total del material.}
\begin{equation}
\Phi_X^{\rm Park}(x)=\int_{\max(1\,{\rm MeV},E_{\rm th})}^{40\,{\rm MeV}}\frac{\Phi_\gamma(E)}{\sigma_{\rm KN}(E)}\frac{\dd\sigma_{\rm prod}(E,x)}{\dd x}\dd E,\quad
\Phi_\gamma(E)=0.58\times10^{18}\frac{P}{\mathrm{MW}}e^{-E/(0.91\,{\rm MeV})}.
\end{equation}
\both{The kernel enforces the additional angular support. Figure 1 uses 1 GW and the three original masses. Open points are vector coordinates extracted from the supplied Park PDF in v2, reused as source data. The original calculation is not reproduced in full; changing the source cutoff or fitting an overall factor is not a verified remedy. The gamma spectrum is an empirical fit above 0.2 MeV, while Park's stated integration cutoff is 1 MeV.}
{El núcleo impone además el soporte angular. La figura 1 usa 1 GW y las tres masas originales. Los puntos abiertos son coordenadas vectoriales extraídas del PDF de Park en v2 y reutilizadas como datos fuente. No se reproduce el cálculo original completo; cambiar el corte del espectro o ajustar un factor global no es una solución verificada. El espectro gamma es un ajuste empírico por encima de 0.2 MeV; el corte de integración indicado por Park es 1 MeV.}
\begin{figure}[htbp]\centering\includegraphics[width=\linewidth]{code/darkphoton_v3/output/figure1_comparison.pdf}\caption{\both{Recalculation versus Park Figure 1. Solid curves are computed; markers are extracted source data.}{Recálculo frente a la figura 1 de Park. Las curvas continuas son calculadas; los puntos son datos extraídos de la fuente.}}\end{figure}
\section{\both{Detection by the inverse process}{Detección por el proceso inverso}}
\both{Time reversal gives the same invariant squared amplitude. For an unpolarized incident massive vector the average over three polarizations adds the factor two thirds relative to production. For incident energy $E$, momentum $p_X$, and outgoing photon energy $w$,}
{La inversión temporal da la misma amplitud invariante al cuadrado. Para un vector masivo incidente no polarizado, el promedio sobre tres polarizaciones introduce el factor dos tercios respecto de producción. Para energía incidente $E$, momento $p_X$ y energía saliente del fotón $w$,}
\begin{align}
p_X&=\sqrt{E^2-M^2},\quad s=m^2+M^2+2mE,\quad a=M^2+2mE,\quad b=-2mw,\\
w_\pm&=\frac{(s-m^2)(E+m\pm p_X)}{2s},\\
\frac{\dd\sigma_{Xe\to\gamma e}^{\rm unpol}}{\dd w}&=\frac23\frac{\epsilon^2e^4(\hbar c)^2 F}{32\pi m p_X^2}.
\end{align}
\both{The inverse massless limit is two thirds of Klein--Nishina. This is appropriate for the Park convention. It is not the factor used for oscillation-produced transverse states. Full containment gives $E_{\rm vis}=E$ after combining electron and photon energy. A real detector needs a response function; replacing it by full containment is an explicit approximation.}
{El límite inverso sin masa es dos tercios de Klein--Nishina. Corresponde a la convención de Park, no a estados transversales producidos por oscilaciones. La contención completa da $E_{\rm vis}=E$ al sumar las energías del electrón y del fotón. Un detector real requiere una función de respuesta; la contención completa es una aproximación explícita.}
\section{\both{Medium production and conversion}{Producción y conversión en el medio}}
\begin{align}
m_\gamma^2&=4\pi\alpha n_e/m,\\
(\Delta m^2)^2&=(M^2-m_\gamma^2)^2+2\epsilon^2M^2(M^2+m_\gamma^2),\\
P_{\gamma\to X}&=\frac{\epsilon^2M^4}{(\Delta m_R^2)^2+E^2\Gamma_R^2},\quad \Gamma_R=\frac{\hbar c}{L_{\rm abs}},\\
P_{X\to\gamma}&=\frac{\epsilon^2M^4}{(\Delta m_D^2)^2},\qquad
\Phi_X^{\rm med}(E)=\Phi_\gamma(E)P_{\gamma\to X}(E).
\end{align}
\both{These are Danilov Eqs. (5), (6), and (9), not all-orders exact probabilities. The earlier implementation incorrectly used 4 instead of 2 in the splitting and treated the detector expression as valid at resonance. Equation (9) explicitly neglects absorption away from resonance. We omit a conservative ten-percent mass window around either plasma mass. A finite-size, absorbing, inhomogeneous detector calculation is needed within it.}
{Son las ecuaciones (5), (6) y (9) de Danilov, no probabilidades exactas a todos los órdenes. La implementación anterior usaba incorrectamente 4 en lugar de 2 en la separación y aplicaba la expresión del detector en resonancia. La ecuación (9) desprecia explícitamente la absorción fuera de resonancia. Omitimos una ventana conservadora del diez por ciento en masa alrededor de cada masa de plasma. Dentro de ella hace falta un cálculo de detector finito, absorbente y no homogéneo.}
\both{We set both plasma masses to 20 eV and the reactor absorption length to 10 cm, following the paper's illustrative choices. At masses well above the plasma scale both probabilities approach the squared mixing. At much smaller masses their product scales as $\epsilon^4(M/m_\gamma)^8$ when damping is negligible. The direct oscillation source preserves photon energy; the Compton-convolved Park source redistributes it. They are alternative source prescriptions and are not multiplied together.}
{Fijamos ambas masas de plasma en 20 eV y la longitud de absorción del reactor en 10 cm, siguiendo las elecciones ilustrativas del artículo. Muy por encima de la escala de plasma ambas probabilidades tienden a la mezcla al cuadrado. Muy por debajo, su producto escala como $\epsilon^4(M/m_\gamma)^8$ si el amortiguamiento es despreciable. La fuente directa de oscilaciones conserva la energía del fotón; la convolución Compton de Park la redistribuye. Son prescripciones de fuente alternativas y no se multiplican entre sí.}
\begin{figure}[htbp]\centering\includegraphics[width=.95\linewidth]{code/darkphoton_v3/output/figure1_medium.pdf}\caption{\both{Medium source at the configured reactor power. The normalization is evaluated at $\epsilon=10^{-6}$ and divided by its square.}{Fuente en el medio a la potencia configurada. La normalización se evalúa en $\epsilon=10^{-6}$ y se divide por su cuadrado.}}\end{figure}
\section{\both{Event rate, statistics, and comparisons}{Tasa de eventos, estadística y comparaciones}}
\begin{align}
N_e&=M_{\rm det}[{\rm kg}]\,10^3N_A Z/A,\\
N_{\rm med}&=\frac{N_eT\eta}{4\pi R^2}\int_{3\,{\rm MeV}}^{8\,{\rm MeV}}\Phi_\gamma(E)P_{\gamma\to X}(E)P_{X\to\gamma}(E)\sigma_{\rm KN}(E)\dd E,\\
N_{\rm med}(\epsilon_{95},M)&=195.7,\quad 0<\epsilon_{95}\le0.01.
\end{align}
\both{Inputs are 2.9 GW, 28 m, 187 kg of CsI, 160 live days, $Z=108$ and $A=259.8099219$ g/mol. The measured excess quoted by Danilov is $30.7\pm100.6$ events. Their one-sided cap is 195.7 selected events. Their assumed anti-Compton suppression corresponds to $\eta=1/6$, approximately 1174 preselection events. It is an estimate, not a measured dark-photon efficiency. We also show unit efficiency to expose its effect. The earlier report used the selected-event cap with unit efficiency only and therefore omitted this correction.}
{Las entradas son 2.9 GW, 28 m, 187 kg de CsI, 160 días de exposición, $Z=108$ y $A=259.8099219$ g/mol. El exceso que cita Danilov es $30.7\pm100.6$ eventos. Su límite unilateral es de 195.7 eventos seleccionados. La supresión anti-Compton que supone equivale a $\eta=1/6$, aproximadamente 1174 eventos antes de selección. Es una estimación, no una eficiencia medida para fotones oscuros. Mostramos también eficiencia unitaria. El informe anterior usaba sólo el límite de eventos seleccionados con eficiencia unitaria y omitía esta corrección.}
\both{The Park-convention recalculation uses the convolved source, massive inverse cross section, unit efficiency and $1.96(100.6)$ events. Park labels this 95\% CL; for a centered Gaussian it corresponds to a 97.5\% one-sided cap. It must not be silently equated with the Danilov one-sided convention. The red line is Park's quoted $2.1\times10^{-5}$ bound, not a newly digitized data set. The black line is extracted directly from Danilov's vector PDF. The independent medium calculation is not tuned to either line. Differences remain because the published normalization, source prescription and detector acceptance are not fully reproduced.}
{El recálculo con la convención de Park usa la fuente convolucionada, la sección inversa masiva, eficiencia unitaria y $1.96(100.6)$ eventos. Park lo llama 95\% CL; para una gaussiana centrada corresponde a un límite unilateral de 97.5\%. No debe identificarse sin explicación con la convención unilateral de Danilov. La línea roja es el límite $2.1\times10^{-5}$ citado por Park, no un conjunto de datos digitalizado de nuevo. La negra se extrae directamente del PDF vectorial de Danilov. El cálculo independiente en el medio no se ajusta a ninguna línea. Persisten diferencias porque no se reproducen completamente la normalización publicada, la prescripción de fuente y la aceptación del detector.}
\begin{figure}[htbp]\centering\includegraphics[width=\linewidth]{code/darkphoton_v3/output/figure2_texono_medium.pdf}\caption{\both{TEXONO-only comparison. Gaps near 20 eV mark the unclaimed resonance region. Medium curves end at 10 keV; the exact massive Park-convention calculation extends to 1 MeV.}{Comparación sólo para TEXONO. Los huecos cerca de 20 eV indican la resonancia sin límite calculado. Las curvas en el medio terminan en 10 keV; el cálculo masivo con la convención de Park llega a 1 MeV.}}\end{figure}
\section{\both{Validity and actual verification}{Validez y verificación realizada}}
\both{The production kernel has exact two-body massive kinematics within tree-level free-electron QED. The medium detector model uses relativistic transverse modes, $M\ll m_e,E$, and is reported only up to 10 keV for a 3--8 MeV ROI. Plasma masses and absorption are constant inputs. Material attenuation, atomic binding, rescattering, escape of photons, finite geometry, and energy-dependent selection are not simulated. Thus the medium limit is conditional on this response and source model; it is not a validated detector-level TEXONO exclusion.}
{El núcleo de producción tiene cinemática masiva exacta de dos cuerpos dentro de QED a nivel árbol para electrones libres. El modelo de detector en el medio usa modos transversales relativistas, $M\ll m_e,E$, y se presenta sólo hasta 10 keV para una ROI de 3--8 MeV. Las masas de plasma y la absorción son entradas constantes. No se simulan atenuación material, ligadura atómica, redispersión, escape de fotones, geometría finita ni selección dependiente de la energía. El límite en el medio está condicionado a este modelo de respuesta y fuente; no es una exclusión TEXONO validada a nivel de detector.}
\both{The executed checks include independent matrix traces and Ward identities, the Reference 15 total, analytic integration, both massless scattering limits, the corrected splitting coefficient, root closure, doubled quadrature resolution, and rejection of invalid domains. Every plotted curve is checked for finite in-frame points and its presence in the canvas at export. The exported images are inspected visually. The original full Park spectrum and Danilov numerical curve are not reproduced by the independent prediction; their differences are visible.}
{Los controles ejecutados incluyen trazas matriciales independientes e identidades de Ward, la sección total de la referencia 15, integración analítica, ambos límites de dispersión sin masa, el coeficiente corregido, cierre de la raíz, doble resolución de cuadratura y rechazo de dominios inválidos. Cada curva se comprueba por sus puntos finitos dentro de los ejes y su presencia en el lienzo al exportar. Las imágenes se inspeccionan visualmente. La predicción independiente no reproduce el espectro completo de Park ni la curva numérica de Danilov; las diferencias quedan visibles.}
\section{\both{Decay survival and the requested sub-MeV sensitivity}{Supervivencia al decaimiento y sensibilidad sub-MeV}}
\both{This execution adds propagation to the checked kernels. It does not derive the electron loop. Izaguirre, Krnjaic and Pospelov, Phys. Rev. D 92, 095014 (2015), Eqs. (9)--(11), distinguish absorption after survival from decays inside a detector. The source is available locally as \texttt{papers/park\_reproduction/reference20.pdf}. Their electron-pair decay channel is closed for the masses considered here. We assume no lighter hidden particles and use the three-photon approximation quoted in Park Eq. (4):}
{Esta ejecución agrega propagación a los núcleos verificados. No deriva el lazo electrónico. Izaguirre, Krnjaic y Pospelov, Phys. Rev. D 92, 095014 (2015), Ecs. (9)--(11), distinguen absorción después de propagación de decaimientos dentro del detector. La fuente local es \texttt{papers/park\_reproduction/reference20.pdf}. El canal de decaimiento a pares electrónicos está cerrado para las masas consideradas. Suponemos que no existen partículas ocultas más livianas y usamos la aproximación a tres fotones citada en la Ec. (4) de Park:}
\begin{align}
\Gamma_{3\gamma}^{\rm approx}&=2.16\times10^{-16}e^4\epsilon^2\frac{M^9}{m_e^8},\\
\ell(E,M,\epsilon)&=\frac{\hbar c\sqrt{E^2-M^2}}{M\Gamma_{3\gamma}},\qquad S(E)=\exp[-R/\ell(E)],\\
N_{\rm vac}&=\frac{N_eT\eta}{4\pi R^2}\int_{E_{\min}}^{E_{\max}}\Phi_X^{\rm Park}(E)S(E)\sigma_{Xe\to\gamma e}(E)\,\dd E.
\end{align}
\both{Both source and scattering include their squared mixing. For the medium model, insert the same survival factor in the earlier medium event integral; do not multiply the two alternative production models. The exponent uses the exact momentum instead of setting the speed to unity. For a detector chord $L$, a decay signal would contain $S(1-e^{-L/\ell})$ and a separate acceptance; we do not count those events as Compton absorption. A detector geometry and decay selection are not provided.}
{La fuente y la dispersión incluyen cada una su mezcla al cuadrado. Para el modelo en el medio, se inserta la misma supervivencia en la integral de eventos anterior; no se multiplican los dos modelos alternativos de producción. El exponente usa el momento exacto en vez de fijar la velocidad a la unidad. Para una cuerda $L$ del detector, una señal de decaimiento llevaría $S(1-e^{-L/\ell})$ y una aceptación propia; no contamos esos eventos como absorción Compton. No se proporciona geometría ni selección de decaimientos.}

\both{The three-photon width is a low-mass expansion, controlled for $M\ll m_e$. Extending it to the upper endpoint below 1 MeV is a diagnostic, not a checked exact-loop prediction. The table records the largest propagation exponent at the vacuum boundary. This small approximate exponent is not an uncertainty bound on the exact width. Invisible decays, if allowed by extra dark particles, would require a different model.}
{El ancho a tres fotones es una expansión de masa baja, controlada para $M\ll m_e$. Extenderla al extremo superior por debajo de 1 MeV es un diagnóstico, no una predicción verificada del lazo exacto. La tabla registra el mayor exponente de propagación en el borde de vacío. Este exponente aproximado pequeño no acota la incertidumbre del ancho exacto. Los decaimientos invisibles, si hay partículas oscuras adicionales que los permitan, requieren otro modelo.}
\input{work/reactor_sensitivity/output/summary.tex}
\begin{figure}[htbp]\centering
\includegraphics[width=\linewidth]{work/reactor_sensitivity/output/conditional_limits.pdf}
\caption{\both{Same selected-event cap and efficiency in both hypotheses. The vacuum curve illustrates the massive Compton model; it is not a valid low-mass material exclusion. The medium curve is restricted to its relativistic validity range and excludes the resonance window.}{Mismo límite de eventos seleccionados y eficiencia en ambas hipótesis. La curva de vacío ilustra el modelo Compton masivo; no es una exclusión material válida a masa baja. La curva en el medio está restringida a su régimen relativista y excluye la ventana resonante.}}
\end{figure}

\section{\both{Reproduction status and reuse}{Estado de reproducción y reutilización}}
\both{The only PDF in \texttt{papers\_init} is Park. Its Figure 1 is recalculated and compared with its vector paths: the heavier-mass spectra disagree, so the full figure is not reproduced. Figure 2 combines TEXONO, NEOS and external astronomical, cosmological and laboratory exclusions. We calculate the TEXONO component, which also disagrees with the printed bound. The original external data and the private NEOS background analysis are not supplied; the complete Figure 2 is not reproduced. Showing the published page or digitizing it would not change this status.}
{El único PDF en \texttt{papers\_init} es Park. Su figura 1 se recalcula y compara con sus trazos vectoriales: los espectros de mayor masa discrepan y la figura completa no está reproducida. La figura 2 combina TEXONO, NEOS y exclusiones astronómicas, cosmológicas y de laboratorio. Calculamos la componente TEXONO, que también discrepa con el límite impreso. No se aportan los datos externos originales ni el análisis privado del fondo de NEOS; la figura 2 completa no está reproducida. Mostrar la página publicada o digitalizarla no cambia ese estado.}

\both{Run \texttt{bash work/reactor\_sensitivity/run.sh} from the project root. Edit \texttt{config.conf} for power, baseline, detector mass, exposure, composition, ROI, efficiency, medium parameters and the event cap. Supply a separately justified \texttt{N95\_events} for a different likelihood, or use the documented Gaussian excess/error option. Background is represented by the excess and its uncertainty, not by a simulated spectrum. The current sub-MeV scan requires an ROI lower edge above 1 MeV. General lower-energy detectors need atomic and response modeling before reuse.}
{Ejecutar \texttt{bash work/reactor\_sensitivity/run.sh} desde la raíz. Editar \texttt{config.conf} para potencia, distancia, masa del detector, exposición, composición, ROI, eficiencia, parámetros del medio y límite de eventos. Proporcionar \texttt{N95\_events} justificado por separado para otra verosimilitud, o usar la opción gaussiana documentada de exceso/error. El fondo está representado por el exceso y su incertidumbre; no por un espectro simulado. Este barrido sub-MeV exige que la ROI comience por encima de 1 MeV. Para otros detectores de menor energía se necesita modelar ligadura atómica y respuesta.}

\both{New C++ functions implement decay propagation, cached event integration and the first upward crossing of the event cap. ROOT supplies quadrature and rendering; existing project code supplies the checked trace and source convolution. No new microscopic amplitude is claimed. The executed checks cover scaling, zero-distance and zero-mixing survival, analytic root recovery, absent roots, quadrature convergence and event-cap closure. Unchecked: the exact near-threshold loop, reactor transport, detector response, polarization transport in a Compton-produced beam, raw TEXONO likelihood and the full published exclusion compilation.}
{Las funciones nuevas en C++ implementan propagación, integración de eventos con pesos precalculados y el primer cruce ascendente del límite de eventos. ROOT aporta cuadratura y gráficos; el código previo aporta la traza verificada y la convolución de la fuente. No se afirma una amplitud microscópica nueva. Los controles ejecutados cubren escalas, supervivencia a distancia y mezcla nulas, recuperación de raíces analíticas, raíces ausentes, convergencia de cuadratura y cierre del límite de eventos. Sin verificar: lazo exacto cerca del umbral, transporte en el reactor, respuesta del detector, transporte de polarización del haz Compton, verosimilitud original de TEXONO y compilación completa de exclusiones.}

\begin{thebibliography}{9}
\bibitem{park} H. K. Park, Phys. Rev. Lett. 119, 081801 (2017), local source \texttt{papers/darkphoton\_v3/park2017.pdf}, Eqs. (1)--(7), Figs. 1--2.
\bibitem{danilov} M. Danilov, S. Demidov and D. Gorbunov, arXiv:1804.10777 (2019), local source \texttt{papers/darkphoton\_v3/danilov2019.pdf}, Eqs. (5), (6), (9), Fig. 1 and TEXONO discussion.
\bibitem{ref15} P. Gondolo and G. Raffelt, Phys. Rev. D 79, 107301 (2009), vector appendix, local source \texttt{papers/darkphoton\_v3/reference15.pdf}.
\end{thebibliography}\end{document}
